% !TEX root = ../thesis.tex
% ==============================================================================
%  CHAPTER 3: LITERATURE REVIEW (فصل سوم: پیشینه پژوهش) (chapters/chapter3_literature.tex)
% ==============================================================================
%  This chapter provides a comprehensive review, taxonomy, and critical
%  comparative survey of existing paradigms and analyzes their theoretical limitations.
% ==============================================================================

\chapter{پیشینه‌ی پژوهش}\label{chap:chapter3_literature}

فصل سوم پایان‌نامه به بررسی نظام‌مند پژوهش‌های پیشین، طبقه‌بندی رهیافت‌های موجود و تبیین شکاف تحقیقاتی اختصاص دارد. هدف این فصل، ارائه‌ی تصویری دقیق و مستند از پیشرفت‌های علمی در حوزه‌ی مسئله است تا مشخص شود چه راهکارهایی تاکنون پیشنهاد شده، هر یک با چه محدودیت‌ها یا مفروضاتی مواجه بوده‌اند، و چرا رویکرد جدید ارائه‌شده در این پایان‌نامه ضرورت دارد. در تدوین پیشینه باید از خلاصه‌نویسی جداگانه‌ی مقالات پرهیز کرد و متون علمی را بر پایه‌ی مفاهیم، شاخص‌ها و روش‌ها به‌صورت تحلیلی مورد بررسی قرار داد.

در این قالب، فصل پیشینه افزون بر ارائه‌ی راهنمای تدوین مرور نوشتار، شیوه‌های استناددهی علمی به انواع مراجع (مجلات، کنفرانس‌ها، کتاب‌ها و پایان‌نامه‌ها) را به نمایش می‌گذارد.


\section{طبقه‌بندی و دسته‌بندی رویکردهای موجود}\label{sec:3-1}

برای ارائه‌ی ساختاریافته‌ی پیشینه، نخستین گام تعریف یک دسته‌بندی منطقی بر اساس مفروضات و اصول روش‌شناسی است. در مسائل علمی و مهندسی، رویکردها معمولاً در چند دسته‌ی کلی جای می‌گیرند:
\begin{itemize}
    \item \textbf{رویکردهای تحلیلی و کلاسیک}: روش‌هایی که بر مبنای مدل‌های صوری و فرمول‌بندی‌های قطعی استوارند~\cite{sample_book, sample_inproceedings}. این روش‌ها شفافیت تحلیلی بالایی دارند، اما تعمیم آن‌ها به سیستم‌های پیچیده یا با ابعاد بالا غالباً با دشواری روبه‌رو است.
    \vspace{-0.25cm}
    \item \textbf{رویکردهای تخمینی و بهینه‌سازی عددی}: روش‌هایی که از الگوریتم‌های تکراری و محاسبات عددی برای همگرایی به پاسخ بهینه بهره می‌برند~\cite{sample_article, sample_inbook, sample_unpublished}.
    \vspace{-0.25cm}
    \item \textbf{رویکردهای داده‌محور و یادگیرنده}: مدل‌هایی که با استفاده از شواهد تجربی و داده‌های جمع‌آوری‌شده رفتار سیستم را تقریب می‌زنند~\cite{sample_incollection, sample_techreport, sample_manual, sample_booklet}.
    \vspace{-0.25cm}
    \item \textbf{رویکردهای ترکیبی}: سامانه‌هایی که از ادغام مدل‌های تحلیلی و روش‌های محاسباتی برای بهره‌گیری هم‌زمان از دقت و مقیاس‌پذیری استفاده می‌کنند~\cite{sample_proceedings, sample_misc, sample_phdthesis, sample_mastersthesis}.
\end{itemize}


\section{تحلیل انتقادی و مقایسه رویکردها}\label{sec:3-2}

پس از طبقه‌بندی مراجع، نقد و بررسی نقاط قوت و ضعف هر دسته از اهمیت بالایی برخوردار است. پژوهشگر باید به‌روشنی مشخص سازد که هر رهیافت تحت چه شرایطی کارایی دارد و در چه شرایطی با ناکارآمدی یا نقایص ساختاری روبه‌رو می‌شود.

عمده‌ترین چالش‌هایی که معمولاً در بررسی تحلیلی روش‌ها شناسایی می‌شوند شامل پیچیدگی محاسباتی بالا، حساسیت شدید به پارامترهای تنظیمی، محدودیت در مقیاس‌پذیری برای داده‌های حجیم و عدم پایداری در حضور نویز یا عدم‌قطعیت‌های محیطی است. ارائه‌ی این تحلیل‌های مقایسه‌ای در قالب متن، تسلط نویسنده بر مرزهای دانش را نمایان می‌سازد.


\section{تبیین شکاف تحقیقاتی و نتیجه‌گیری فصل}\label{sec:3-3}

بخش پایانی فصل پیشینه باید به بیان صریح «شکاف پژوهشی» بپردازد. شکاف پژوهشی همان پرسش یا نیاز پاسخ‌داده‌نشده‌ای است که کارآمدی روش‌های پیشین را با بن‌بست روبه‌رو کرده و انگیزه و دلیل اصلی طراحی متدولوژی پیشنهادی را توجیه می‌نماید.

با جمع‌بندی ادبیات تحقیق مشخص می‌شود که علی‌رغم پژوهش‌های گسترده، دستیابی به روشی که هم‌زمان پیچیدگی محاسباتی اندک، پایداری بالا و سهولت پیاده‌سازی را فراهم سازد، همچنان نیازمند مطالعه است. این ضرورت، زمینه‌ی ورود به فصل چهارم و ارائه‌ی روش‌شناسی پیشنهادی را فراهم می‌سازد.
